\section{Objetivos} \label{objetivos}

El objetivo principal de este trabajo es desarrollar una plataforma funcional y validada de
comunicación con periféricos serie sobre el MPSoC Zynq UltraScale+ ZCU102, que sirva
como base para el OBC del proyecto LINCE. Para alcanzarlo, se plantearon los siguientes
objetivos específicos:

\vspace{0.2cm}
\noindent\textit{Conviene advertir que los seis objetivos no se plantearon simultáneamente. El
encargo inicial comprendía únicamente los objetivos 1, 2 y 5 ---el transceptor en la lógica
programable, su driver para RTEMS y las herramientas de automatización del flujo---. El objetivo 3
se incorporó cuando el diseño de una tarjeta de expansión propia, abordado por iniciativa propia
para poder validar los catorce canales, llevó al consorcio a asignar también las dos tarjetas
oficiales del proyecto. El objetivo 4, la comparativa de arquitecturas de transporte, surgió del
propio desarrollo: la arquitectura de partida resultó incapaz de sostener catorce canales
simultáneos, y determinar su sustituta se convirtió en la aportación de mayor recorrido del
trabajo.}
\vspace{0.2cm}

\begin{enumerate}
    \item \textbf{Diseñar un IP core de transceptor serie configurable en VHDL} para la lógica
    programable del MPSoC, capaz de operar sobre los estándares RS422 y RS485, con
    parámetros configurables en tiempo de ejecución: baudrate, paridad, número de bits de
    datos, bits de parada y orden de bit.

    \item \textbf{Desarrollar un driver de software completo en C para RTEMS} que abstraiga el
    acceso al hardware y permita a la aplicación gestionar hasta 14 instancias de transceptor
    simultáneas mediante una API orientada a interrupciones, sin recurrir a \textit{polling} de CPU.

    \item \textbf{Diseñar y fabricar las tarjetas de expansión de prueba} requeridas para validar el
    sistema frente a los subsistemas CDHS (gestión de datos y calentadores) y AOCS (control
    de actitud y órbita) del satélite, integrando interfaces RS422/RS485, CAN, SPI, SpaceWire
    y PWM.

    \item \textbf{Determinar la arquitectura de transporte de datos adecuada entre el PS y la PL}
    para catorce canales serie simultáneos. Para ello se implementarán tres arquitecturas
    completas ---AXI GPIO, catorce controladores AXI DMA y un AXI MCDMA multicanal con un puente
    diseñado en VHDL--- y se compararán con un mismo banco de pruebas en términos de ocupación de
    la FPGA, carga de interrupciones sobre la CPU, latencia y \textit{throughput}.

    \item \textbf{Desarrollar herramientas de automatización del flujo de desarrollo}, incluyendo
    scripts TCL para la generación automática del diseño en Vivado y scripts de despliegue de
    imágenes, con el fin de reducir los tiempos de iteración.

    \item \textbf{Integrar y validar el sistema completo}, verificando la comunicación entre la
    ZCU102 y las tarjetas de expansión a través de los distintos buses implementados.
\end{enumerate}
